\chapter*{Resumen}
\addcontentsline{toc}{chapter}{Resumen}

La presente investigación desarrolla y valida una metodología basada en
criterios técnicos y regulatorios para la selección de tecnologías de
interfaz de usuario en el contexto de transformación digital de instituciones
bancarias, centrándose en el análisis comparativo de Angular, React y Vue.js.
En un sector donde la seguridad, el rendimiento y el cumplimiento normativo
resultan críticos, la elección tecnológica representa una decisión
estratégica con impacto directo en la competitividad y en la exposición
regulatoria de la institución.

Un factor adicional motiva y delimita el alcance de esta investigación: la
consolidación de herramientas de generación de código asistida por
Inteligencia Artificial (GitHub Copilot, Cursor y asistentes equivalentes), que
reduce progresivamente la brecha de productividad y de curva de aprendizaje
asociada al talento disponible para cada framework --un factor
tradicionalmente organizacional-- y desplaza el eje decisivo de la selección
hacia las propiedades técnicas intrínsecas del framework y su compatibilidad
regulatoria, que la IA no puede sustituir. Por ello, el modelo propuesto se
estructura en dos dimensiones --técnica y regulatoria-- con diez criterios de
evaluación, dos de ellos emergentes: la calidad y verificabilidad del código
generado por IA (dimensión técnica) y la gobernanza regulatoria del uso de
herramientas de IA generativa en el desarrollo de software bancario
(dimensión regulatoria).

La metodología propuesta consta de cuatro fases integradas: análisis de
requisitos institucionales, evaluación técnica (incluida la calidad del código
asistido por IA), evaluación regulatoria y de cumplimiento, y un proceso
estructurado de decisión final. Para validar este enfoque, se implementó un
caso de estudio en una institución bancaria mediana en proceso de
modernización de sus canales digitales, aplicando sistemáticamente los
criterios y herramientas definidos en la metodología, y contrastando los
resultados con y sin la incorporación de los criterios asociados a la IA
generativa.

Los resultados demuestran que, bajo el modelo técnico-regulatorio, Angular
obtiene la puntuación ponderada más alta ($S=8.84$), seguido de React
($S=7.83$) y Vue.js ($S=7.56$), gracias a su arquitectura tipada y opinionada,
que favorece tanto la seguridad y mantenibilidad del código como la
trazabilidad exigida por auditorías regulatorias y, adicionalmente, produce
código asistido por IA más consistente y verificable. El análisis de
sensibilidad confirma que esta preferencia es robusta: a diferencia de un
modelo que pondera también factores organizacionales --donde un escenario
extremo podía favorecer a React--, el modelo técnico-regulatorio propuesto
mantiene a Angular como opción óptima en todos los escenarios plausibles
para un entorno bancario regulado.

Esta investigación contribuye al campo con un marco metodológico adaptado
específicamente a las necesidades del sector bancario en la era de la IA
generativa, proporcionando instrumentos de evaluación y decisión que pueden
ser adoptados por instituciones financieras enfrentando desafíos similares de
transformación digital. Se concluye que la irrupción de los asistentes de
código basados en IA no elimina la necesidad de un proceso de selección
tecnológica riguroso, sino que reconfigura sus criterios decisivos hacia lo
técnico y lo regulatorio.

\vspace{0.5cm}
\noindent \textbf{Palabras clave:} Transformación digital bancaria,
tecnologías de interfaz de usuario, Angular, React, Vue, criterios técnicos y
regulatorios, generación de código asistida por Inteligencia Artificial,
metodología de selección tecnológica, aplicaciones financieras.
